\par\Needspace{6\baselineskip}\medskip
\SourceLeaf{140}\pdfbookmark[2]{100.}{article-100}
\begin{center}100.\end{center}
In solving equation 13 or $13^*$, we shall first consider separately the case in which $z$ is small enough for $Z$ to be expressed by a rapidly converging series in powers of $z$. We have
\[
\begin{gathered}
(1+2z)\BareRoot{(z+zz)}=z^{\frac12}+\tfrac52z^{\frac32}+\tfrac78z^{\frac52}-\cdots,\\[.5em]
\log(\BareRoot{(1+z)}+\BareRoot z)=z^{\frac12}-\tfrac16z^{\frac32}+\tfrac3{40}z^{\frac52}-\cdots,
\end{gathered}
\]
\noindent Thus the numerator of $Z=\tfrac83z^{\frac32}+\tfrac45z^{\frac52}-\cdots$ and the denominator $=2z^{\frac32}+3z^{\frac52}+\cdots$, whence $Z=\tfrac43-\tfrac85z+\cdots$. To find the rule governing the series, we differentiate the equation
\[
2(z+zz)^{\frac32}Z=(1+2z)\BareRoot{(z+zz)}-\log(\BareRoot{(1+z)}+\BareRoot z),
\]
\noindent which after simplification gives
\[
2(z+zz)^{\frac32}\frac{\mathrm dZ}{\mathrm dz}+3Z(1+2z)\BareRoot{(z+zz)}=4\BareRoot{(z+zz)}
\]
\noindent or
\[
(2z+2zz)\frac{\mathrm dZ}{\mathrm dz}=4-(3+6z)Z,
\]
\noindent from which, by reasoning similar to that in Art. 90, we obtain
\[
\begin{aligned}
Z={}&\tfrac43-\frac{4\mathpunct{.}6}{3\mathpunct{.}5}z+\frac{4\mathpunct{.}6\mathpunct{.}8}{3\mathpunct{.}5\mathpunct{.}7}zz\\[.5em]
&-\frac{4\mathpunct{.}6\mathpunct{.}\overset{\mkern2mu .}{8}\mathpunct{.}10}{3\mathpunct{.}5\mathpunct{.}7\mathpunct{.}9}z^3+\frac{4\mathpunct{.}6\mathpunct{.}8\mathpunct{.}10\mathpunct{.}12}{3\mathpunct{.}5\mathpunct{.}7\mathpunct{.}9\mathpunct{.}11}z^4-\text{etc.}
\end{aligned}
\]
\noindent It is thus clear that $Z$ depends on $-z$ in exactly the same way as $X$ depended on $x$ above for the ellipse. So if we set
\[
Z=\frac1{\tfrac34+\tfrac9{10}(z+\zeta)}\mathpunct{\raisebox{.45ex}{,}}
\]
\noindent then $\zeta$ is determined by $-z$ in the same way as $\xi$ was determined above by $x$, so that we have
\[
\mathrm{14)}\qquad\zeta=\cfrac{\tfrac2{35}zz}{1-\tfrac2{35}z+\cfrac{\tfrac{40}{63}z}{1+\cfrac{\tfrac4{99}z}{1+\cfrac{\tfrac{70}{143}z}{1+\text{etc.}}}}}
\]
\noindent or
\[
\zeta=\cfrac{\tfrac2{35}zz}{1+\tfrac{18}{35}z+\cfrac{\tfrac4{63}z}{1+\cfrac{\tfrac{40}{99}z}{1+\cfrac{\tfrac{18}{143}z}{1+\text{etc.}}}}}
\]
\noindent\SourceLeaf{141}The values of $\zeta$ were calculated in this way from $z=0$ to $z=0{,}3$ in steps of one thousandth. The third column of Table III gives these values.
